\documentclass[french]{article}
\usepackage{verbatim}
\usepackage{amsmath,amsfonts,amssymb}
\usepackage{pgfplots}
\usepackage[utf8]{inputenc}
\usepackage[french]{babel}
\usepackage{pgf,tikz,tkz-tab}
\usepackage{mathrsfs}
\usetikzlibrary{arrows}
\usepackage{pstricks}
\pgfplotsset{compat=newest}
\usepackage{geometry}
\geometry{hmargin=0.7cm,vmargin=2cm}


\begin{document}
\twocolumn
\begin{center}
\underline{\textbf{Fonctions de référence}}
\end{center}

\section{La fonction carré}

\section*{Exercice 1}
Calculer, lorsque c'est possible, les antécédents des réels suivants par la fonction carré.
\begin{enumerate}
\item 1
\item -16
\item 95
\item 25
\end{enumerate}
%\section*{Exercice 2}
%Soit $f$ la fonction carré. Pour chacune des propositions suivantes indiquer si elle est vraie ou fausse. Justifier.
%\begin{enumerate}
%\item Tous les nombres réels ont exactement une image par $f$.
%\item Il existe un nombre réel qui n'a pas d'antécédent par $f$.
%\item Tous les nombres réels ont, au plus, un antécédent par $f$.
%\item Il existe au moins un nombe réel qui a deux antécédent par $f$.
%\end{enumerate}

\section*{Exercice 2}
Résoudre algébriquement dans $\mathbb{R}$ les équations suivantes.
\begin{enumerate}
\item $x^2=16$
\item  $x^3=-3$
\item  $x^2=15$
\end{enumerate}


%Comparer, sans calcul, les réels suivants. 
%\begin{enumerate}
%\item $(-1,22)^2$ et $(-1,21)^2$
%\item $(\pi -1)^2$ et $(\pi +1)^2$
%\item $(\sqrt{2}-4)^2$  et $(\sqrt{2}-3)^2$
%\end{enumerate}

\section*{Exercice 3}
Résoudre graphiquement dans $\mathbb{R}$ les inéquations suivantes.
\begin{enumerate}
\item $x^2 \leq 9$
\item $x^2 \geq 3$ 
\item $x^2 > 0 $
\item $x^2 \leq -2$
\end{enumerate}

%\section*{Exercice 5}
%Déterminer le sens de variation des fonctions suivantes.
%\begin{enumerate}
%\item $f$ définie sur $\mathbb{R}$ par $f(x)=\dfrac{1}{2}x^2$.
%\item $g$ définie sur $\mathbb{R}$ par $g(x)=7x^2$.
%\item $h$ définie sur $\mathbb{R}$ par $h(x)=-2x^2+7$.
%\end{enumerate}



\section{La fonction inverse}
\section*{Exercice 4}
Soit $g(x)=\dfrac{1}{x}$, déterminer graphiquement la valeur approchée des images suivantes :
\begin{enumerate}
\item $g(4)$
\item $g(0,25)$
\item $g(-3,2)$
\item $g(\dfrac{3}{5})$
\item $g(-\dfrac{7}{5})$ 

Déterminer graphiquement la valeur approchée d'un antécédent \textit{(s'il existe)} dans les cas suivants:

\item $g(x)=4$
\item $g(x)=-0,44$
\item $g(x)=\dfrac{2}{5}$
\end{enumerate}

\section*{Exercice 5}
Résoudre graphiquement dans $\mathbb{R}$ les inéquations suivantes.
\begin{enumerate}
\item $-2 \leq \dfrac{1}{x} \leq 0$ 
\item $1 \leq \dfrac{1}{x} \leq 2$ 
\item $0 \leq \dfrac{1}{x} \leq 3$ 
\item $-2 \leq \dfrac{1}{x} \leq 0$ 
\item $4>\dfrac{1}{x} \geq -1$
\end{enumerate}

\section*{Exercice 6}
Soit $x \in \mathbb{R}^*$, résoudre algébriquement.
\begin{enumerate}
\item $\dfrac{1}{x}=2$
\item $\dfrac{1}{x}=-5$
\item $\dfrac{3}{x}=2$
\item $\dfrac{1}{x}=\dfrac{5}{3}$
\item $\dfrac{1}{x}+3=2$
\item $\dfrac{2}{x}=x$
\item $\dfrac{7}{x}+2=5$
\item $\dfrac{2}{x}=x$
\end{enumerate}

%\section*{Exercice 7}
%Comparer, sans calcul, les réels suivants. 
%\begin{enumerate}
%\item $\dfrac{1}{2,401}$ et $\dfrac{1}{1,202}$
%\item $\dfrac{1}{-3,1}$ et $\dfrac{1}{-3,2}$
%\item $\dfrac{1}{\pi}$  et $\dfrac{1}{3,14}$
%\end{enumerate}



\section*{Exercice 7}
Déterminer le sens de variation des fonctions suivantes et dresser le tableau de variations correspondant.
\begin{enumerate}
\item $f$ définie sur $\mathbb{R}^*$ par $f(x)=\dfrac{1}{x}+3$.
\item $g$ définie sur $\mathbb{R}^*$ par $g(x)=-\dfrac{2}{x}$
\end{enumerate}

\section{La fonction cube et racine carrée}
\section*{Exercice 8}
Calculer les images des nombres suivants par la fonction cube.

\begin{enumerate}
\item $\sqrt{2}$
\item $-\dfrac{3}{4}$
\item $10^3$
\item $\dfrac{\sqrt{5}}{3}$
\end{enumerate}

\section*{Exercice 9}

Comparer, sans calcul, les réels suivants.
\begin{enumerate}
\item $\sqrt{\dfrac{1}{2}}$ et $\sqrt{\dfrac{1}{3}}$.
\item $\sqrt{2,001}$ et $ \sqrt{2,01}$.
\item  $\sqrt{\pi}$ et $\sqrt{3}$.
\end{enumerate}

\section*{Exercice 10}
Résoudre dans $\mathbb{R}_+$ les équations suivantes.
\begin{enumerate}
\item $\sqrt{x} = 2$ 
\item $\sqrt{x} = 3$
\item $-\sqrt{x} = -\dfrac{3}{4}$ 
\end{enumerate}
\end{document}
